\chapter{Exercicis proposats}
\label{cap:exercicis}
\begin{enumerate}
\item Calcula les operacions següents amb complexos:

\begin{enumerate}
\item[a)]
\[
\frac{(2-3i)(1-i)}{(-2+i)(-1-3i)}-\frac{i(3+i)}{2-3i}
\]

\item[b)]
\[
i^{327}-\frac{1}{i^{26}}+\frac{i^{125}}{1-i}
\]

\item[c)]
\[
\left( \frac{i^{7}-1}{i^{29}+1}\right) ^{4}
\]
\end{enumerate}

\begin{indicacio}
(a) $\frac{16}{65}-\frac{41}{65}i$; (b) $\frac{1}{2}-\frac{1}{2
}i$; (c) $1$
\end{indicacio}

\item Resol en $\mathbb{C}$ les equacions següents amb la incògnita $z$:

\begin{enumerate}
\item[a)]
\[
\frac{z}{1+2i}+3+2i=2+5i
\]

\item[b)]
\[
(4+3i)z+7-10i=2+3i+(1-2i)z
\]
\end{enumerate}

\begin{indicacio}
(a) $z=-7+i$; (b) $z=\frac{25}{17}+\frac{32}{17}i$
\end{indicacio}

\item Resol en $\mathbb{C}$ els següents sistemes:

\begin{enumerate}
\item[a)]
\[
\left\{
\begin{array}{l}
(1-i)u+iv=1+3i \\
(2+i)u+(1+i)v=1-3i
\end{array}
\right.
\]

\item[b)]
\[
\left\{
\begin{array}{l}
(3+i)u+(3-i)v=-4 \\
(2+3i)u-(2-3i)v=4
\end{array}
\right.
\]
\end{enumerate}

\begin{indicacio}
(a)
\[
\left\{
\begin{array}{l}
u=-\frac{21}{13}-\frac{1}{13}i \\
v=\frac{19}{13}-\frac{35}{13}i
\end{array}
\right.
\]
(b)
\[
\left\{
\begin{array}{l}
u=\frac{2}{9}+\frac{4}{9}i \\
v=-\frac{10}{9}-\frac{8}{9}i
\end{array}
\right.
\]
\end{indicacio}

\item Calcula la part real i la part imaginària del complex
\[
z=i+\frac{(1-i)^{2}}{3+i}-\frac{(1-i)^{3}}{(2+i)^{2}}
\]

\begin{indicacio}
$\func{Re}(z)=\frac{9}{25}$ i $\func{Im}(z)=\frac{8}{25}$
\end{indicacio}

\item Troba els nombres reals $a$ i $b$ perquè
\[
(a+bi)(2-3i)+(a-bi)(1-i)=2+3i
\]

\begin{indicacio}
$a=-\frac{4}{11},b=\frac{17}{11}$
\end{indicacio}

\item Determina $a\in \mathbb{R}$ perquè $(1+ai)(2-3i)$ sigui un nombre
real.

\begin{indicacio}
$a=\frac{3}{2}$
\end{indicacio}

\item Determina $a\in \mathbb{R}$ perquè $(1-3i)(a-i)-(1+2i)(1-ai)$ sigui un
complex imaginari pur.

\begin{indicacio}
$a=-4$
\end{indicacio}

\item Escriu una equació de segon grau amb coeficients reals de
manera que una de les seves solucions sigui $3-2i$.

\begin{indicacio}
$x^{2}-6x+13=0$
\end{indicacio}

\item Escriu una equació de segon grau que tingui com a solucions $
3+2i$ i $4+i$.

\begin{indicacio}
$x^{2}-(7+3i)x+10+11i=0$
\end{indicacio}

\item Expressa en forma polar els complexos següents:
\[
\begin{array}{llllllll}
a) & -2 &  & b) & -3i &  & c) & i^{1237} \\
&  &  &  &  &  &  &  \\
d) & -2-2\sqrt{3}i &  & e) & \cos 35
^{\circ}
+i\cos 55
^{\circ}
&  & f) & 4(-\cos 15
^{\circ}
+i\sin 15
^{\circ}
)
\end{array}
\]

\begin{indicacio}
(a) $2_{\pi }$; (b) $3_{\frac{3\pi }{2}}$; (c) $1_{\frac{\pi }{
2}}$; (d) $4_{\frac{4\pi }{3}}$; (e) $1_{35
^{\circ}
}$; (f) $4_{165
^{\circ}
}$
\end{indicacio}

\item Si $z$ és un complex unitari d’argument $\alpha $, escriu en
forma trigonomètrica els complexos
\[
\begin{array}{lllllllllll}
a) & z &  & b) & z^{-1} &  & c) & z^{3} &  & d) & z^{2}+\frac{1}{z^{2}}
\end{array}
\]

\begin{indicacio}
(a) $\cos \alpha +i\sin \alpha $; (b) $\cos \alpha -i\sin
\alpha $; (c) $\cos 3\alpha +i\sin 3\alpha $; (d) $2\cos 2\alpha
$
\end{indicacio}

\item Mitjançant la fórmula de Moivre, expressa $\sin 4\alpha $ i $\cos
4\alpha $ en función de $\sin \alpha $ i $\cos \alpha $.

\begin{indicacio}
$\sin 4\alpha =\cos ^{4}\alpha -6\cos ^{2}\alpha \sin
^{2}\alpha +\sin ^{4}\alpha $ i $\cos 4\alpha =4\cos ^{3}\alpha
\sin \alpha -4\cos \alpha \sin ^{3}\alpha $
\end{indicacio}

\item Per què valores de $n\in N$, $(1+i)^{n}$ és imaginari.

\begin{indicacio}
$n=4k-2$, sent $k\in \mathbb{N}$
\end{indicacio}

\item Calcula el mòdul i l’argument dels complexos següents:

\begin{enumerate}
\item[a)]
\[
z_{1}=\frac{(1+i)^{5}}{(1-i)^{3}}
\]

\item[b)]
\[
z_{2}=3\left( \cos \frac{3\pi }{2}+i\sin \frac{3\pi }{2}\right) ^{10}
\]
\end{enumerate}

\begin{indicacio}
(a) $\left\vert z_{1}\right\vert =2$ i $\arg (z_{1})=0$; (b) $
\left\vert z_{2}\right\vert =3$ i $\arg (z_{2})=\pi $
\end{indicacio}

\item Calcula la suma dels cent primers sumands de la suma:
\[
(-1+2i)+(1+i)+3+(5-i)+(7-2i)+\cdots
\]

\begin{indicacio}
$9800+4800i$
\end{indicacio}

\item Calcula les operacions següents amb complexos i expressa el
resultat en forma binòmica:

\begin{enumerate}
\item[a)]
\[
\left( \frac{\sqrt{3}}{2}-\frac{1}{2}i\right) ^{21}
\]

\item[b)]
\[
\left( \frac{3}{\sqrt{3_{30
^{\circ}
}}}\right) ^{6}
\]

\item[c)]
\[
\left[ \frac{\sqrt{2}\left( \cos \frac{\pi }{4}+i\sin \frac{\pi }{4}\right)
}{2\left( \cos \frac{\pi }{6}+i\sin \frac{\pi }{3}\right) }\right] ^{12}
\]
\end{enumerate}

\begin{indicacio}
(a) $i$; (b) $-6\sqrt{3}-6\sqrt{3}i$; (c) $\frac{1}{729}$
\end{indicacio}

\item Escriu en forma trigonomètrica l’expressió següent
\[
\frac{(\cos 5x+i\sin 5x)^{6}\cdot (\cos 3x-i\sin 3x)^{7}}{(\cos x+i\sin
x)^{4}\cdot (\cos 2x-i\sin 2x)^{2}}
\]

\begin{indicacio}
$\cos 9x+i\sin 9x$
\end{indicacio}

\item Calcula el valor de $x$ perquè el producte $1_{30
^{\circ}
}\cdot 5_{x}$ sigui: (a) un nombre real positiu; (b) un nombre real
negatiu; (c) un imaginari pur positiu; (d) un imaginari pur negatiu
imaginari pur negatiu.

\begin{indicacio}
(a) $x=-30
^{\circ}
+k\cdot 360
^{\circ}
$; (b) $x=150
^{\circ}
+k\cdot 360
^{\circ}
$; (c) $x=60
^{\circ}
+k\cdot 360
^{\circ}
$; (d) $x=-60
^{\circ}
+k\cdot 360
^{\circ}
$, sent $k\in \mathbb{Z}$
\end{indicacio}

\item Troba dos nombres complexos sabent que el seu quocient és $3$, els seus arguments sumen $\frac{\pi }{3}$ i els seus mòduls sumen $8$.

\begin{indicacio}
$3(\sqrt{3}+i)$ i $\sqrt{3}+i$
\end{indicacio}

\item Calcula
\[
\sqrt[4]{16\left( \cos 80
^{\circ}
+i\sin 80
^{\circ}
\right) }
\]

\begin{indicacio}
\[
\sqrt[4]{16_{80
^{\circ}
}}=\left\{
\begin{array}{l}
2_{20
^{\circ}
} \\
2_{110
^{\circ}
} \\
2_{200
^{\circ}
} \\
2_{290
^{\circ}
}
\end{array}
\right.
\]
\end{indicacio}

\item Calcula les arrels cinquenes de la unitat imaginària i representa-les gràficament.

\begin{indicacio}Calcula les cinc arrels com a arrels cinquenes de $i=1_{90^{\circ}}$ i representa-les separades per angles de $72^{\circ}$.\end{indicacio}

\item Resol en $\mathbb{C}$ les equacions següents:

\begin{enumerate}
\item[a)]
\[
z^{3}-i=1
\]

\item[b)]
\[
z^{3}-2z^{2}+9z-18=0
\]

\item[c)]
\[
z^{6}-2z^{3}+2=0
\]

\item[d)]
\[
z^{4}-z^{3}+z^{2}-z+1=0
\]

\item[e)]
\[
z^{2}-(10i-5)z-14i-48=0
\]
\end{enumerate}

\begin{indicacio}
(a) $z_{1}=\left( \sqrt[6]{2}\right) _{\frac{\pi }{12}
},z_{2}=\left( \sqrt[6]{2}\right) _{\frac{9\pi }{12}},z_{3}=\left( \sqrt[6]{2
}\right) _{\frac{17\pi }{12}}$; (b) $z_{1}=2,z_{2}=3i,z_{3}=-3i$; (c) $
z_{1}=\left( \sqrt[6]{2}\right) _{15
^{\circ}
},z_{2}=\left( \sqrt[6]{2}\right) _{135
^{\circ}
},z_{3}=\left( \sqrt[6]{2}\right) _{255
^{\circ}
},z_{4}=\left( \sqrt[6]{2}\right) _{-45
^{\circ}
},z_{5}=\left( \sqrt[6]{2}\right) _{105
^{\circ}
},z_{6}=\left( \sqrt[6]{2}\right) _{225
^{\circ}
}$; (d) $z_{1}=1_{36
^{\circ}
},z_{2}=1_{108
^{\circ}
},z_{3}=1_{252
^{\circ}
},z_{4}=1_{324
^{\circ}
}$; (e) $z_{1}=-8+6i,z_{2}=3+4i$
\end{indicacio}

\item Troba dos complexos conjugats tals que la seva diferència és $6i$ i el seu
quocient és imaginari pur.

\begin{indicacio}
$3+3i,3-3i$ o bé $-3+3i,-3-3i$
\end{indicacio}

\item Troba dos complexos tals que un és imaginari pur, el mòdul
de l’altre és $\frac{\sqrt{2}}{2}$ i la seva suma és igual al seu producte.

\begin{indicacio}
$i,\frac{1}{2}-\frac{1}{2}i$ o bé $-i,\frac{1}{2}+\frac{1}{2}
i$
\end{indicacio}

\item Troba dos complexos tals que el seu producte és $4i$ i el seu quocient és
\thinspace $2_{\frac{\pi }{6}}$.

\begin{indicacio}
$\sqrt{2}+\sqrt{6}i$ i
$\frac{\sqrt{6}}{2}+\frac{\sqrt{2}}{2}i$
\end{indicacio}

\item Troba els vèrtexs d’un hexàgon regular de centre l’origen,
si se sap que un d’ells és $(0,1)\,$.

\begin{indicacio}
$(0,1),\left( -\frac{\sqrt{3}}{2},\frac{1}{2}\right) ,\left( -
\frac{\sqrt{3}}{2},-\frac{1}{2}\right) ,(0,-1),\left( \frac{\sqrt{3}}{2},-
\frac{1}{2}\right) ,\left( \frac{\sqrt{3}}{2},\frac{1}{2}\right)
$
\end{indicacio}

\item D’un triangle equilàter es coneixen un vèrtex $(0,2)$ i
el baricentre $(0,0)$. Troba els altres dos vèrtexs.

\begin{indicacio}
$(-\sqrt{3},-1)$ i $(\sqrt{3},-1)$
\end{indicacio}

\item Un rombe els quals costats formen un angle de 45
$^{\circ}$
té dos vèrtexs consecutius en $(0,0)$ i $(1,1)$. Troba els altres
dos vèrtexs si se sap que estan en el primer quadrant.

\begin{indicacio}
Dos solucions: $(1,1+\sqrt{2}),(0,\sqrt{2})$ i $(1+\sqrt{2}
,1),(\sqrt{2},0)$
\end{indicacio}

\item Donats els complexos $u=1+i$ i $v=4+2i$, troba el complex $z$ que
els afixos de $u,v,z$ són els vèrtexs d’un triangle equilàter.

\begin{indicacio}
Dos solucions: $z=\frac{1}{2}(5-\sqrt{3})+\frac{1}{2}(3\sqrt{3
}+3)i$ i $z=\frac{1}{2}(5+\sqrt{3})+\frac{1}{2}(3-3\sqrt{3})i$
\end{indicacio}

\item Troba els vèrtexs d’un quadrat del qual se sap que el centre és
$(3,3)$ i el punt mitjà d’un costat és $(4,1)\,$.

\begin{indicacio}
$(2,0),(6,2),(4,6),(0,4)$
\end{indicacio}

\item Quina transformació geomètrica representa
l’aplicació de $\mathbb{C}$ en $\mathbb{C}$ definida per $z^{\prime
}=2+iz$?

\begin{indicacio}
És un gir d’angle recte i centre punt $(1,1)$
\[
z-z_{C}=(z-z_{C})\cdot 1_{\frac{\pi }{2}}
\]
amb $C=(1,1)$.
\end{indicacio}

\end{enumerate}

